% Pursuing Stacks (cotes 134-2 to 134-8): archive cote 134-7, site page 65 (PDF page 66), typescript page 507.
% typescript. Draft compared: ps6.tex:3172-3233.
% AI-transcribed by Claude Opus 5.5 (claude-opus-5-5, adaptive thinking, effort high), Message Batches API, batch ps4 (msgbatch_012Aoc1NuskATPsyBeKSYK7J), 2026-11-12; not yet checked by eye.
\scanpage{134-7/65}
\mspage{507}
compatible with the structures of $k_o$-algebras as well as composition
operations,  and satisfying moreover two conditions for
\[ F(x+y) \ , \quad \text{resp. } F(xy) \]
in terms of two diagonal maps
\[ \diagram{\begin{tikzcd} \hand{\Omega}_o \ar[r,shift left,"\hand{\Delta_a}"] \ar[r,shift right,"\hand{\Delta_m}"'] & \hand{\Omega}_1 = \hand{\Omega}_o\otimes_{k_o}\hand{\Omega}_o \end{tikzcd}} \]
(which may be described in terms of the composition structure on $\hand{\Omega}_o$ ,
as expressing the compositions  $Fo(G'+G")$  resp.  $Fo(G'G")$),  and moreo\xout{x}-%
ver one trivial condition for $F(\hand{\lambda}))$  when $\hand{\lambda}$ in $k$ comes from $k_o$ , \xout{\illegible{1}}
namely compatibility of the map $k_o \to k$  with the operations of $\hand{\Omega}_o$ on
both $k_o$ and $k$. (NB T\unsure{h}e operation of $\hand{\Omega}_o$ on $k_o$ is defined by
\[ F(\hand{\lambda}) = Fo\hand{\lambda} \ , \]
where $k_o$ is identified to a subring of $\hand{\Omega}_o$ .)

I did'nt pursue much further these ponderings,  just one digression
among many in the main line of investigation !  I also read through
the preprint of Davi\over{c}{d} W. Jones on Poly-T-complexes,  which Ronnie Brown
(\inserted{acting as}David Jones' supervisor) had sent me a while ago.  There he develops
a notion of polyhedral cells,  with a view of using these instead of
simplices or cubes for doing combinatorial homotopy theory\unsure{.|,} \xout{in a spirit}
As I had pondered a little along this direction (cf. sections 91\hand{,} topic
\hand{(}8),  and section 93),  I was h\unsure{o}ping that some of the perplexities I
had been meeting would be solved in David Jone\unsure{s}\xout{\unsure{k}}\hand{'} notes  -  for instance
that there would be handy criteria for a category $M$ made up with such
polyhedral cells to be a weak test category,  namely that objects of $M^{\wedge}$
may be used as models for homotopy types;  also\hand{,}  that the``standard''
\xout{\unsure{simple}} chain complex constructed in section 93 is indeed an ``abelianiza-%
tor'' for $M$,  i.e. may be used for computing homology of objects of $M^{\wedge}$.
David Jones' emphasis however is a rather different one  -  he seems
mainly interested in generalizing the theory of ``thin'' structures
of M.K. Dakin \over{\unsure{t}}{f}rom the simplicial to the more general polyhedral set-up,
and prove a corresponding equivalence of categories.  Thus my perplexi-%
ties remain - they are admittedly rather marginal in the main line of
thought,  and I doubt I'll stop to try and solve them.

\margin{65}
